% \begin{Details}

% \begin{Aside}
% Beam finite element formulations have, over several decades, proliferated into a large and fragmented collection of special-case implementations.
% Production codes typically offer distinct element types for thin-walled sections, for asymmetric sections, for large displacements, for warping torsion, and so on---each developed and validated in relative isolation.
% An improvement to one formulation does not automatically benefit others; a correction to the handling of, say, shear-flexure interaction in a geometrically linear element may have no bearing on the corresponding geometrically nonlinear element, even within the same software package.
% The result is a maintenance burden that grows combinatorially with the phenomena one wishes to capture, and a validation effort that must be repeated for each combination.

% This fragmentation also obscures deficiencies.
% Certain widely-used formulations for geometrically nonlinear analysis violate strain objectivity: objective strain measures, which should be invariant under superposed rigid motions, are corrupted by the interpolation procedure \citep{crisfield1999objectivity}. 
% The errors are small for moderate rotations and vanish at converged equilibrium states, so they are often tolerated in practice.
% But tolerating such errors becomes untenable when simulations must be differentiated---as required for sensitivity analysis, gradient-based inference, or optimization---because the spurious dependence on rigid motion contaminates the computed gradients.
% Fixing the problem within existing software architectures would require rewriting element-level code, with attendant re-validation, for each affected element type.
% \end{Aside}

\subsection{Motivation}\label{sec:part-i-objectives}



% A family of models are developed for simulating slender members exhibiting geometric and material nonlinearities. 
Slender members are typically represented by dimensionally reduced \citep{kantorovich1958approximate} models of three dimensional continuua, known as beam theories, which simplify the partial differential equations of elastostatics to produce an ordinary differential equation in one spatial variable.
When discretized by the finite element method, these models offer a substantial advantage in computational demand, but several complications arise in the face of geometric and material nonlinear response:
\begin{description}
\item[Geometric Nonlinearity.]
Dimensional reduction of the geometrically nonlinear continuum equations generally leads to a problem posed over a constrained solution space that lacks a vectorial structure. 
Such problems lie outside the domain of classical finite element methods \citep{strang1973analysis}, which can fail to preserve important features of the constrained model, like strain objectivity and path independence \citep{jelenić1999geometrically}. 
Formulations that restore objectivity exist, but they sacrifice the symmetric tangent stiffness that accelerates iterative solution.

\item[Material Nonlinearity.]
Inelastic material response exhibits a complicated dependency on the three dimensional state of stress and strain in a body. 

Fiber section models \citep{taucer1991fiber} reconstruct a \emph{uniaxial} strain field by leveraging the plane section hypothesis of Euler-Bernoulli and Timoshenko beam theories. 
This approach yields acceptable accuracy when a uniaxial material model can sufficiently capture the inelastic material response. 
This is often the case for slender members with simple cross sections that are not expected to experience torsion. 

It is possible to reconstruct a \emph{multiaxial} strain field from the fields of a beam theory by leveraging the plane section hypothesis, but this fails to account for cross-sectional warping and tends to over-predict strength and stiffness. 

A more accurate approach for reconstructing a multiaxial strain state introduces auxiliary kinematic fields in the formulation of the beam theory that model this warping \citep{reissner1952nonuniform,klinkel2003anisotropic,lecorvec2012nonlinear}. 
When no limit is placed on the number of warping fields, it is possible to represent a continuum to any desired accuracy. 
With only one auxiliary warping field, \cite{armero2021modeling} has demonstrated that such models furnish stress fields which violate equilibrium and may produce inaccurate results. 
However, excessive auxiliary warping fields can be computationally prohibitive. 
Furthermore such models are often cumbersome, and it is preferable to forgo auxiliary warping fields altogether whenever possible. 

% In the elastic theory of beams these effects are tacitly represented by the Saint~Venant torsion and shear-correction cross sectional constants. 

% \item[Physical Transparency.] [TODO: Setup for how certain formulations of shear and rotations are required to model certain boundary conditions and loadings, and the tradeoffs that arise.]
\end{description}

Beam theories which are capable of representing multiaxial nonlinear material response fall into two broad categories:
\begin{itemize}
    \item Type I theories employ only the classical deflection and rotation fields as unknowns and furnish finite element implementations with six degrees of freedom per node.
    \item Type II theories include at least one auxiliary field to model cross sectional warping and furnish finite element implementations with more than six degrees of freedom per node.
\end{itemize}

Objective: Multiaxial material response in beams.



\subsection{Literature}


\subsubsection{Type I}

When cross sectional warping is expected to vary uniformly over the body, auxiliary warping fields can be eliminated by associating various warping modes with classical beam deformations like curvature \(\bm{\kappa}\). 
The phenomenon of torsional warping is particularly amendable to this approach and has been successfully incorporated into beam finite elements by several authors \citep{gruttmann1998geometrical}, but warping induced by transverse shear introduces several difficulties:
\begin{enumerate}
    \item Warping is not naturally proportional to a classical beam deformation field
    \item The shear warping shapes arising in Saint-Venant's solution do not exhibit the same orthogonality properties as the torsion warping function.
\end{enumerate}
In 2D, (2) is avoided and a satisfactory remedy for (1) may be obtained by appealing to classical shear-correction arguments.
An example of such a treatment is the investigation by \cite{simo1982bifurcation} which employed Cowper's construction \citep{cowper1966shear} to produce a two-dimensional beam theory of elastomeric bearings. However, as elaborated in \cref{sec:timo}, this approach cannot be made variationally consistent in 3D. 
A similar two dimensional approach was taken by \cite{saritas2006mixed} to formulate a beam theory of shear-yielding structural steel members.

In the case of pure bending, it is standard that inelastic elements maintain exact consistency with Saint-Venant's pure bending solution \citep[Art.87]{love1944treatise}, which finds bending resistance in proportion to the elastic modulus \(E\) and the second-area-tensor of the cross section \(\bm{I}\) \citep{spacone1992beam,ibrahimbegovic2023geometrically}. Similar consistency in torsion and elongation is also standard \citep{gruttmann1998geometrical}. 
There does not appear to be a formulation for shear that accomplishes this achievement. 

\subsubsection{Type II}
When cross sectional warping is \emph{nonuniform}, auxiliary warping fields become necessary. 
However, it is still desirable to minimize the number of such fields. 
These additional degrees 
of freedom either inflate the global system of equations—incurring a cost 
that scales as the cube of the DOF count for direct linear solvers—or must 
be condensed out at the element level during each Newton iteration.

\cite{lecorvec2012nonlinear} put forth a Type II formulation where warping patterns were manually constructed for various thin-walled shapes using Lagrange interpolation. 
These interpolation schemes had to be constructed uniquely for a given section geometry, and were not able to reproduce classical Timoshenko beam solutions in the elastic range. 
Furthermore, due to the general nature of the Lagrange interpolation basis, several warping fields were required to accurately reproduce classical solutions. 

\subsection{Scope and Methodology}

An assumed strain framework is developed that enhances section constitutive models for standard beam elements and supports integrating inelastic multiaxial material response over the cross section. 


Beam theories are typically derived for finite element analysis by one of two approaches:
\begin{itemize}
\item In a \emph{constrained continuum} approach \citep{antman1995nonlinear}, the field equations of continuum mechanics are stated weakly in the form of a functional\footnote{Typical functionals include virtual work, Hu-Washizu, and Hellinger-Reissner.} 
and the displacement field is postulated.
\end{itemize}

% The developments are restricted to the following setting :
% \begin{itemize}
% \item Prismatic, initially straight beams with arbitrary cross-section geometry. 
% \item Quasi-static response.
% \item Small to moderate material strain.
% % \item Initially isotropic.
% \end{itemize}


\begin{remark}\label{rem:inelastic-consistency}
Although the objective of this work is enhanced simulation in the nonlinear regime, the methodology of \cref{sec:timo} and \cref{sec:nonuniform} which targets linear elastic consistency is naturally a necessary feature of an inelastic model, and follows a standard motif:
\begin{itemize}
    \item The enhanced strain fields employed for inelastic plane finite elements \citep{simo1990class} are derived from those exhibited in elastic bending \citep{wilson1990use}.
    \item The plane-section strain profile employed by standard fiber beam elements \citep{spacone1992beam} ensures agreement with the elastic solution for bending  \citep[Art. 87]{love1944treatise}.
    \item Saint Venant's linear elastic torsional warping function is employed in the geometrically nonlinear beam of \citep{simo1991geometricallyexact} and the inelastic beam of \citep{gruttmann2000theory}.
\end{itemize}
\end{remark}

\begin{remark}
By admitting an arbitrary number of auxiliary warping modes, the general beam theory outlined in \cref{sec:genwarp} generalizes the scope of classical plane-section beam theories, whose correspondence with continuum mechanics is only achieved asymptotically with the slenderness of a member \citep{parker1979asymptotic,marigo2006hierarchy,mielke1988saintvenant}. %cimetiere1988asymptotic
Each auxiliary warping field enriches the function space of the approximate model in a manner similar to mesh refinement of a continuum finite element model.
With few judiciously selected auxiliary warping fields, the beam theory of \cref{sec:genwarp} can achieve the same accuracy of a continuum finite element model with a significant reduction in computational demand. 
\end{remark}


\subsection{Outline}

% \Cref{sec:genwarp} develops equilibrium equations for a family of beam theories that account for cross-sectional warping in the geometrically exact context.
% Beam theories are classified into four types according to how warping amplitudes are related to the classical beam strains. 
% The Type~3 formulation, in which warping amplitudes are algebraically constrained to the beam strains, is of particular interest as it requires only six degrees of freedom.
\begin{description}
\item[Preliminaries.]
\Cref{sec:genwarp} introduces the general mathematical setting. 
Notation is introduced and the theory of geometrically exact Cosserat beams with auxiliary warping is recalled. 
The configuration manifold of the dimensionally reduced model is postulated and general governing equilibrium equations are derived by imposing objectivity. 
The distinction between uniform and nonuniform warping is introduced and discussed. 


\item[Enhanced Strain Framework.]
\Cref{sec:gensec} presents a novel enhanced strain principle for formulating constitutive equations in the setting of \cref{sec:genwarp}, and presents an algorithm for its implementation. 
This general principle is specialized in \Cref{sec:timo} and \Cref{sec:nonuniform} for uniform and nonuniform warping, respectively. 
In both cases, the enhanced shapes are derived to ensure continuum-consistent response in the elastic regime (\cref{rem:inelastic-consistency}). 

\item[Enhanced Uniform Warping.]
\Cref{sec:timo} develops enhanced shapes for situations where warping is expected to be uniform. 
The objective is to achieve agreement with Saint Venant's linear elastic continuum solutions for loads in the \(\mathrm{P}_{\rm I\!I}\) class (\cref{sec:sv}). 
Example~\ref{sec:frame-0012} demonstrates continuum consistency of the resulting section model using standard beam elements without auxiliary warping degrees of freedom.

\item[Enhanced Nonuniform Warping.]
\Cref{sec:nonuniform} develops enhanced shapes for situations where non-uniform warping is expected to be significant.
Continuum solutions are not generally available for these problems, and consistency is instead achieved by requiring pointwise equilibrium in the elastic regime. 
\end{description}




%through the reduced Type~3 models of \cref{sec:genwarp}. 
% The investigation gives rise to the 

% \end{Details}

